\section{Conclusion}\label{sec:conclusion}

This thesis set out to relax one of the two assumptions behind semiparametric estimates of the effects of
monetary policy. Its central observation holds up: the Federal Reserve's own briefing documents show that some
actions are not on the table at some meetings, and the Committee never chose an action its staff had not
proposed. Treating the Bluebook menu as a covariate that dictates where the propensity score is zero gives a
decision model that fits markedly better at no cost in parameters, and a sharp characterisation of what the data
can and cannot say about the average effect of policy.

What the data can say is limited. Without restrictions on the menus that exclude the action or the status quo,
the average effect is not signed. With effects assumed similar across menus, the estimates are precise but have
the wrong sign for output and unemployment, even with the Greenbook forecasts in the propensity score. The
forecasts do remove the price puzzle for rate increases. The contribution of the menu is therefore to isolate the
problem. Overlap can be handled; unconfoundedness is what fails.

Section~\ref{sec:discussion} sets out the natural next steps: a richer, text-based information set inside the
menu-constrained propensity score, a sensitivity analysis for unconfoundedness alongside the one for effect
homogeneity, and the combination of the menu restriction with high-frequency policy surprises.
